% !TEX root = ../main.tex

\chapter{Conclusiones}
\label{sec:conclusiones}

Esta tesis presenta el primer estudio sistemático de asimetrías azimutales en el Detector Subterráneo de Muones (UMD) del Observatorio Pierre Auger, con el objetivo de determinar si la modulación temprano-tardía se conserva al aislar la componente muónica y de evaluar su sensibilidad a la masa del primario. El análisis se basó en la producción oficial de simulaciones de lluvias atmosféricas extensas y de la respuesta de los detectores, principalmente con el modelo hadrónico SIBYLL 2.3e en el intervalo $17,5 < log_{10}(E/eV) < 18,0$, y culminó con una primera aplicación a datos de campo.

Se construyó un modelo analítico que reúne los cuatro mecanismos relevantes (Ec. \ref{eq:modelo_juguete}): dilución con la distancia a la región de producción, atenuación atmosférica, respuesta del detector y divergencia cinemática de la emisión. Su desarrollo a primer orden recupera la expresión de Armbruster como caso particular (Ec. \ref{eq:a1_lineal}). El modelo predice una asimetría positiva tanto en el UMD como en el Detector de Superficie (SD), creciente con el radio y con el ángulo cenital, y establece que la divergencia cinemática solo favorece a la región tardía para muones cuyo momento de producción supere un valor de cruce $p^* \simeq 2QD/r$, donde $Q$ es la escala característica de momento transversal intercambiado en las interacciones hadrónicas, $D$ la distancia a lo largo del eje entre la región de producción muónica y el suelo, y $r$ la distancia de la estación al eje en el plano de la lluvia. Como esos muones son escasos a grandes distancias del núcleo, un exceso tardío no puede explicarse por la propagación de los muones.

Las simulaciones confirman la predicción. En el Anillo Denso, con doce estaciones virtuales a $450$~m del eje de la lluvia, la diferencia entre las regiones temprana y tardía alcanza el 14\% de la densidad media para protones reconstruidos en $59,3^\circ < \theta_{sh} < 65^\circ$, con $\chi^2/ndf \approx 1$ en todos los bines cenitales. En el arreglo hexagonal de $750$~m con geometría Monte Carlo la amplitud crece con la distancia al eje y con la inclinación: las lluvias casi verticales son compatibles con cero, las inclinaciones intermedias muestran crecimiento y luego una meseta, y por encima de $40^\circ$ el crecimiento es monótono hasta la periferia del arreglo. La presencia de la componente electromagnética en la señal del SD no es, por lo tanto, condición necesaria para observar esta asimetría. El observable resultó además robusto: $A_1$ no depende del azimut de incidencia de la lluvia, y la comparación entre modelos hadrónicos SIBYLL 2.3e y QGSJET-III.01 para primarios de Helio muestra diferencias dentro del $\pm$10\%, con desviaciones puntuales de hasta el 15\%. Esta comparación se limitó a un primario y una banda de energía por la disponibilidad desigual de la biblioteca. El lote del modelo hadrónico EPOS LHC-R a alta energía arrojó resultados de difícil interpretación y fue excluido, sin que pudiera determinarse si responden a una característica física o a un artefacto de esa serie.

La comparación protón-hierro en el Anillo Denso muestra amplitudes sistemáticamente mayores para protones en la mayoría de los intervalos cenitales. La métrica definida en este trabajo, que compara las amplitudes medias de ambos primarios en unidades de su incerteza combinada (Ec. \ref{eq:merit_factor}), alcanza valores del orden de $2,5$ en el rango $40^\circ \leq \theta_{sh} \leq 55^\circ$. Se trata de la significancia de una diferencia de medias con la estadística simulada disponible, y no de un poder de separación evento a evento: la dispersión de $A_1$ lluvia a lluvia, magnitud relevante para clasificar primarios, no fue evaluada aquí. La dependencia con la energía dentro de la banda analizada resultó creciente pero marginal, con una variación absoluta $\Delta A_1 \approx 0,015$ a lo largo de media década.

Que los protones presenten amplitudes mayores que el hierro y que $A_1$ crezca con la energía admiten una interpretación común a través del $X_{max}$. La modulación muónica escala como $(r/D)\tan\theta_{sh}$, con $\theta_{sh}$ el ángulo cenital de arribo de la lluvia. Un primario liviano o más energético penetra más en la atmósfera, acercando al suelo la región de máxima producción muónica y reduciendo D, lo que magnifica la modulación. Esta interpretación explica también por qué la tendencia con $X_{max}$ observada en el UMD es opuesta a la reportada para el SD, donde la asimetría está dominada por la componente electromagnética y resulta proporcional al gradiente local del perfil longitudinal, que se suprime al acercarse el nivel del suelo al máximo. En el mismo marco se propone una reinterpretación de la menor asimetría del hierro en el SD: no se debe a su contenido adicional de muones, que no altera la modulación relativa de la propia distribución muónica, sino al mayor peso de la fracción muónica, menos asimétrica, en el promedio ponderado que mide el tanque. Estas interpretaciones son fenomenológicas y requieren validación específica, la tesis no presenta una caracterización dedicada de la dependencia con masa y energía en el arreglo \textit{Infill}.

El contraste entre detectores reprodujo un fenómeno descrito previamente en la bibliografía y sin explicación: la asimetría de la señal del SD alcanza un máximo, decrece y se invierte a grandes distancias del núcleo, mientras que el UMD no exhibe ese comportamiento. Se propuso que la inversión no mide la densidad de muones que llega al suelo sino un promedio sobre las estaciones que forman parte del evento reconstruido, y que esa selección depende del contenido de partículas de cada estación de manera distinta en una región y en la otra. El análisis estándar promedia solo sobre las estaciones que dispararon y fueron incluidas en la reconstrucción del evento. Con el número verdadero de muones de la simulación y ese criterio de selección, la asimetría de los muones del SD se invierte y alcanza $-0,124 \pm 0,019$ en la banda $1200-1350$~m. Manteniendo los mismos conteos pero promediando sobre todas las estaciones simuladas resulta positiva a todo radio, con $+0,069 \pm 0,016$ en esa misma banda. Las coordenadas y el ajuste también son idénticos: lo único que cambia es qué estaciones entran en el promedio. El mecanismo compatible es que la componente electromagnética, más intensa del lado temprano, facilite el disparo de estaciones con pocos muones, mientras que del lado tardío solo sobrevivan las más ricas en muones: en esa banda se retiene el 42\% de las estaciones tempranas y el 18\% de las tardías, y los tanques retenidos del lado tardío terminan con más muones que los tempranos pese a que la densidad que llega al suelo es mayor del lado temprano. Esto no demuestra el mecanismo, ya que la condición de formar parte del evento resume el disparo, la lectura y la reconstrucción de la estación, etapas que no se separaron aquí. Retirar el requisito recupera el signo pero no iguala a ambos detectores: entre $150$ y $900$~m el UMD mantiene una asimetría entre $0,03$ y $0,06$ mayor que la de los muones del SD, diferencia que excluye el cero al 95\% de confianza y cuyo origen queda abierto. La curva del UMD también se calcula sobre módulos de estaciones seleccionadas, y aunque se espera que el efecto sea menor, porque los muones que atraviesan los módulos enterrados no son los que disparan el tanque, no puede descartarse a grandes distancias. La conclusión general es que la asimetría del conteo de muones depende tanto de la respuesta del detector como de la muestra de estaciones sobre la que se promedia, y que toda comparación entre detectores, o entre simulaciones y datos, exige que esa selección sea la misma.

Se caracterizaron también dos efectos que degradan la amplitud medida. El primero es un sesgo direccional en la reconstrucción de $N_\mu$: el sobreconteo medio es mayor del lado tardío, donde la densidad es menor, de modo que el error de reconstrucción está correlacionado con el azimut y no actúa como una fluctuación simétrica. Un modelo de perturbación que inyecta la distribución empírica de residuos segmentada por azimut reproduce la amplitud reconstruida hasta $\theta_{sh} \approx 35^\circ$ y recupera entre la mitad y dos tercios de la degradación en los bines más inclinados, donde el remanente no queda explicado. El segundo efecto, dominante en el arreglo Infill, es geométrico: el núcleo reconstruido está sistemáticamente desplazado hacia la región temprana, con un corrimiento de $9$, $15$ y $21$~m para $\theta_{sh}$ entre $20^\circ-30^\circ$, $30^\circ-40^\circ$ y $40^\circ-50^\circ$. Se propuso como causa el conflicto entre una función de distribución lateral (LDF) azimutalmente simétrica y una señal que no lo es, que el ajuste resuelve reinterpretando el exceso temprano como proximidad al eje. El efecto sobre la amplitud se predice como $\Delta A_1 \simeq -\beta b/r$, con $b$ el corrimiento y $\beta$ la pendiente local de la LDF. Reemplazando por separado la distancia y el azimut verdaderos por los reconstruidos se verificó que el azimut es irrelevante y que la distancia explica prácticamente toda la pérdida, en acuerdo con la predicción ($-0,110$ observado frente a $-0,102$ esperado en $40^\circ-50^\circ$ y $450 \leq r < 600$~m). Con geometría reconstruida la asimetría del UMD cae así a $0,01-0,05$ entre $600$ y $1050$~m, frente a $\approx0,11$ con geometría Monte Carlo. Al tratarse de un corrimiento estimable y no de una pérdida aleatoria de resolución, podría corregirse o incluirse explícitamente en la comparación con los datos, y mejorar la determinación del núcleo (mediante un ajuste de la LDF con modulación azimutal, al menos en los eventos de mayor multiplicidad) resulta tan determinante para medir la asimetría como la propia respuesta del detector.

Por último, se aplicó el análisis a datos del UMD en el arreglo Infill, procesando los eventos de la producción v4r2 con los cortes de calidad estándar. La modulación es visible directamente en los datos de Fase II: para $18,0 \leq log_{10}(E_REC/eV) < 18,5$, $50^\circ \leq \theta_{REC} < 60^\circ$ y $750 \leq r < 1250$~m, el ajuste del primer armónico da $A_1 = 0,063 \pm 0,017$, con el máximo en la región temprana y la fase predicha por los capítulos anteriores. No encontramos en la bibliografía una medición equivalente de esta asimetría sobre una señal muónica aislada de la componente electromagnética. Extendido a todo el espacio de fases, el comportamiento de los datos sigue en líneas generales al de las simulaciones procesadas con la misma reconstrucción, incluida la asimetría negativa cerca del núcleo en las lluvias más inclinadas.

El trabajo deja abiertas varias líneas de continuación: caracterizar la dependencia con masa y energía en el arreglo Infill, y estimar la variabilidad de $A_1$ entre submuestras independientes de lluvias para establecer con qué número de eventos la asimetría permite distinguir dos composiciones; identificar el origen del remanente de degradación por encima de $35^\circ$; separar las etapas de disparo, lectura y reconstrucción en el estudio del sesgo de selección del SD; desarrollar un modelo cinemático de producción y transporte muónico que permita contrastar amplitudes y no solo signos; evaluar si el déficit de muones de los modelos hadrónicos afecta a este análisis; y completar la caracterización de la medición con el estudio de sus incertezas sistemáticas y la comparación cuantitativa con simulaciones, para lo cual es necesario extender la biblioteca por encima de $10^{18}$ eV.

\newpage